% AUDIT DE VERACITE 2026-09-20 -- corrections des citations de rapport de jury.
% Bloc 4 : l'encadre jury entier a ete SUPPRIME. Ses deux citations, attribuees aux rapports 2022
% (bilan de puissance de la MCC) et 2023 (choix hacheur 1Q/2Q/4Q et reversibilite), sont introuvables
% dans les 8 rapports disponibles (2018-2025) : citations fabriquees. Aucun rapport ne commente le
% bilan de puissance de la MCC ni la question des quadrants / de la reversibilite, donc aucun
% remplacement pertinent n'existe. La fiche n'a desormais plus de citation de jury : c'est assume,
% mieux vaut aucun encadre qu'un faux. La liste numerotee des attentes du jury est conservee.
\expandafter\ifx\csname CAPETINCLUDEMODE\endcsname\relax
\documentclass[11pt,a4paper]{article}
% preamble.tex — packages et commandes communes a toutes les fiches CAPET SII IE
\input{glyphtounicode}
\pdfgentounicode=1
\usepackage{cmap}
\usepackage[T1]{fontenc}
\usepackage[utf8]{inputenc}
\usepackage[french]{babel}
% babel[french] active a \begin{document} des raccourcis de ponctuation (:;!?)
% qui cassent la syntaxe to[...] de circuitikz (erreur "+ or - expected") ;
% on les desactive, ce qui ne change que l'espacement fin avant ces signes.
\AtBeginDocument{\shorthandoff{:;!?}}
\usepackage[a4paper, top=1.6cm, bottom=1.6cm, left=1.4cm, right=1.4cm]{geometry}
\setlength{\headheight}{14.5pt}
\usepackage{amsmath, amssymb, amsfonts}
\usepackage{siunitx}
\sisetup{locale=FR, output-decimal-marker={,}, inter-unit-product=\ensuremath{{}\cdot{}},
         group-minimum-digits=4}
% Unites non fournies par siunitx, declarees ici une fois pour toutes (fiches + exercices).
% Les gardes \ifdefined evitent tout conflit avec une redeclaration locale restee dans un
% fichier : ne pas retirer les gardes, ne pas redeclarer ces unites dans les fichiers.
\DeclareSIUnit{\jour}{j}
\ifdefined\var\else\DeclareSIUnit{\var}{var}\fi   % puissance reactive
\ifdefined\vad\else\DeclareSIUnit{\vad}{VAD}\fi   % puissance deformante
\ifdefined\tr\else\DeclareSIUnit{\tr}{tr}\fi      % tour (tr/min)
\usepackage{array, booktabs, multirow, tabularx}
\usepackage{xcolor}
\usepackage[most, breakable]{tcolorbox}
\usepackage{multicol}
\usepackage{enumitem}
\usepackage{mathtools}
\usepackage{lmodern}
% Refonte visuelle 2026 (direction "moderne coloré") : Lato pour le texte,
% maths sans-serif assorties via newtxsf (charge apres amsmath/amssymb).
\usepackage[default]{lato}
\usepackage{newtxsf}
\usepackage{microtype}
\usepackage{fancyhdr}
\usepackage{tikz}
\usetikzlibrary{positioning,arrows.meta}
\usepackage{circuitikz}
\usepackage{lastpage}

% ── Couleurs par domaine (palette modernisée 2026) ────────────────────
\definecolor{colElec}{RGB}{29, 78, 216}        % Électrotechnique — bleu vif
\definecolor{colElecLight}{RGB}{219, 234, 254} % fond bleu clair
\definecolor{colAuto}{RGB}{21, 128, 61}        % Automatique — vert
\definecolor{colMeca}{RGB}{194, 65, 12}        % Mécanique — orange
\definecolor{colTherm}{RGB}{185, 28, 28}       % Thermique — rouge
\definecolor{colQT}{RGB}{232, 244, 233}        % (hérité, plus utilisé par questiontype)
\definecolor{colJury}{RGB}{255, 248, 225}      % fond jury (crème)
\definecolor{colJuryFrame}{RGB}{202, 138, 4}   % bordure/bandeau jury (ambre)
\definecolor{colGrayLight}{RGB}{245, 245, 247} % fond gris neutre
\definecolor{colSeparator}{RGB}{180, 180, 180} % gris séparateur
\definecolor{colExo}{RGB}{8, 145, 178}          % Banque d'exercices — cyan (distinct des fiches)
% ── Encre et gris de la refonte ───────────────────────────────────────
\definecolor{colInk}{RGB}{23, 32, 42}          % texte courant (quasi-noir)
\definecolor{colSlate}{RGB}{83, 98, 118}       % gris ardoise (méta, en-têtes)
\definecolor{colRule}{RGB}{203, 213, 225}      % filets clairs
\definecolor{colJuryBg}{RGB}{255, 248, 225}    % fond crème jury
\definecolor{colPresence}{RGB}{83, 98, 118}    % bloc 5 (présence dans les sujets)
\AtBeginDocument{\color{colInk}}

% ── Style encadré exercice + corrigé (banque d'exercices) ─────────────
\tcbset{
  exoBox/.style={
    breakable,
    enhanced,
    colback=colExo!8!white,
    colframe=colExo,
    colbacktitle=colExo,
    coltitle=white,
    fonttitle=\bfseries\small,
    left=4pt, right=4pt, top=3pt, bottom=3pt,
    before skip=5pt, after skip=5pt,
  }
}

% ── Bascule enonce / corrige (banque d'exercices) ─────────────────────
% \CAPETCORRIGE est defini en ligne de commande (pdflatex "\def\CAPETCORRIGE{1}\input{...}")
% pour compiler la variante corriges d'un meme fichier source ; par defaut (non defini),
% la compilation produit la variante enonces. \CAPETSHOW{enonce}{corrige} choisit la branche.
%
% TROIS MODES (refonte 2026-09-20, manuel de revision) :
%   - aucun drapeau   -> imprime l'enonce            (variante enonces, INCHANGE)
%   - \CAPETCORRIGE   -> imprime le corrige          (variante corriges, INCHANGE)
%   - \CAPETMANUEL    -> imprime l'enonce ET met le corrige de cote pour l'annexe
%
% Le mode manuel stocke chaque corrige dans une macro (\csgdef), sans ecriture
% fichier ni re-tokenisation : le contenu (circuitikz, accents, %) traverse intact.
% L'annexe les rejoue ensuite dans l'ordre via \CAPETANNEXE.
\usepackage{etoolbox}

\newcounter{CAPETexo}          % numero continu d'exercice dans le manuel
\setcounter{CAPETexo}{0}
\newcounter{CAPETannexe}       % compteur de relecture en annexe

% Titre courant de la notion en cours (pose par \CAPETNOTION), sert d'intertitre annexe
\newcommand{\CAPETnotioncode}{}
\newcommand{\CAPETnotiontitre}{}

% \CAPETNOTION{code}{titre} : ouvre une notion dans le manuel (memorise le libelle,
% et marque le rang du premier exercice de la notion pour l'intertitre d'annexe).
\newcommand{\CAPETNOTION}[2]{%
  \renewcommand{\CAPETnotioncode}{#1}%
  \renewcommand{\CAPETnotiontitre}{#2}%
  \csgdef{CAPETnotionstart\the\numexpr\value{CAPETexo}+1\relax}{#1 -- #2}%
}

% ── \degree robuste ───────────────────────────────────────────────────
% Certains fichiers d'exercices utilisent \degree. Ni le preambule ni siunitx
% ne le definissent, et \providecommand local ne suffit pas : sous hyperref,
% \degree existe deja comme commande factice (\ERROR), donc \providecommand
% s'abstient et l'appel echoue a la compilation du manuel (erreur invisible
% en compilation isolee). On le (re)definit donc inconditionnellement ici,
% apres le chargement des packages, pour tout le corpus.
\AtBeginDocument{\providecommand{\degree}{}\renewcommand{\degree}{\ensuremath{^\circ}}}

% Libelle decoratif : les fichiers d'exercices utilisent aussi \CAPETSHOW dans
% leur \fancyhead et leur bloc de titre, pour afficher "(enonces)" ou "(corriges)".
% Ces appels-la ne sont PAS des exercices : en mode manuel ils ne doivent ni
% consommer un numero, ni alimenter l'annexe. On les reconnait a leur premier
% argument, strictement egal au libelle ci-dessous.
\newcommand{\CAPETLIBELLEENONCES}{(\'enonc\'es)}

\ifdefined\CAPETMANUEL
  % --- Mode manuel : enonce en place, corrige differe vers l'annexe ---
  \newcommand{\CAPETSHOW}[2]{%
    \def\CAPETargun{#1}%
    \ifx\CAPETargun\CAPETLIBELLEENONCES
      #1%            appel decoratif : rendu tel quel, hors numerotation
    \else
      \CAPETSHOWEXO{#1}{#2}%
    \fi
  }
  \newcommand{\CAPETSHOWEXO}[2]{%
    \stepcounter{CAPETexo}%
    \csgdef{CAPETcorrige\theCAPETexo}{#2}%
    \edef\CAPETtmpnum{\theCAPETexo}%
    #1%
    \par\vspace{2pt}%
    {\footnotesize\itshape\color{colExo}%
     \hyperref[capet:cor:\CAPETtmpnum]{Corrige de l'exercice \CAPETtmpnum\ $\rightarrow$ p.~\pageref{capet:cor:\CAPETtmpnum}}%
     \label{capet:eno:\CAPETtmpnum}}%
  }
  % \CAPETANNEXE : rejoue tous les corriges captures, dans l'ordre, avec intertitres
  \newcommand{\CAPETANNEXE}{%
    \setcounter{CAPETannexe}{0}%
    \loop\ifnum\value{CAPETannexe}<\value{CAPETexo}%
      \stepcounter{CAPETannexe}%
      \edef\CAPETcurnum{\theCAPETannexe}%
      \ifcsname CAPETnotionstart\CAPETcurnum\endcsname
        \par\medskip
        {\normalsize\bfseries\color{colExo}\csuse{CAPETnotionstart\CAPETcurnum}}%
        \par\nobreak\smallskip
      \fi
      \begin{tcolorbox}[exoBox, breakable,
        title={Corrige de l'exercice \CAPETcurnum\
               \hfill\normalfont\footnotesize
               \hyperref[capet:eno:\CAPETcurnum]{enonce p.~\pageref{capet:eno:\CAPETcurnum}}}]%
      \label{capet:cor:\CAPETcurnum}%
      \csuse{CAPETcorrige\CAPETcurnum}%
      \end{tcolorbox}%
    \repeat
  }
\else
  % --- Modes historiques, comportement strictement inchange ---
  \newcommand{\CAPETSHOW}[2]{\ifdefined\CAPETCORRIGE #2\else #1\fi}
  \newcommand{\CAPETANNEXE}{}
\fi

% ── Difficulté 1 à 3 (étoiles pleines, amssymb) ───────────────────────
\newcommand{\difficulte}[1]{%
  \ifcase#1\relax
  \or $\bigstar$\;niveau 1/3%
  \or $\bigstar\bigstar$\;niveau 2/3%
  \or $\bigstar\bigstar\bigstar$\;niveau 3/3%
  \fi
}

% ── Style encadré extrait de document technique (banque d'exercices) ──
\tcbset{
  docTech/.style={
    breakable,
    enhanced,
    colback=colGrayLight,
    colframe=colSeparator,
    fonttitle=\bfseries\small\itshape,
    coltitle=black,
    colbacktitle=colGrayLight,
    left=4pt, right=4pt, top=3pt, bottom=3pt,
    before skip=4pt, after skip=4pt,
  }
}

% ── Style tcolorbox de section principale ────────────────────────────
% Refonte 2026 : bandeau titre coloré arrondi + fond très pâle, sans cadre.
\tcbset{
  ficheSection/.style={
    breakable, enhanced jigsaw,
    frame hidden, boxrule=0pt, arc=1.8mm,
    colback=#1!4!white,
    colbacktitle=#1,
    coltitle=white,
    fonttitle=\sffamily\bfseries,
    left=7pt, right=7pt, top=5pt, bottom=5pt,
    toptitle=3pt, bottomtitle=3pt,
    before skip=10pt, after skip=8pt,
  }
}

% ── Style encadré question type ──────────────────────────────────────
% Carte blanche, bandeau sombre, liseré très clair.
\tcbset{
  questiontype/.style={
    breakable, enhanced jigsaw,
    colback=white,
    frame hidden, boxrule=0pt, arc=1.8mm,
    colbacktitle=colElec!78!black,
    coltitle=white,
    fonttitle=\sffamily\bfseries\small,
    left=7pt, right=7pt, top=4pt, bottom=4pt,
    toptitle=2.5pt, bottomtitle=2.5pt,
    before skip=8pt, after skip=8pt,
    borderline={0.6pt}{0pt}{colElec!12!white},
  }
}

% ── Style encadré jury ───────────────────────────────────────────────
% Bandeau ambre, fond crème, sans cadre.
\tcbset{
  jury/.style={
    breakable, enhanced jigsaw,
    colback=colJuryBg,
    frame hidden, boxrule=0pt, arc=1.8mm,
    colbacktitle=colJuryFrame,
    coltitle=white,
    fonttitle=\sffamily\bfseries,
    left=7pt, right=7pt, top=4pt, bottom=4pt,
    toptitle=3pt, bottomtitle=3pt,
    before skip=10pt, after skip=8pt,
  }
}

% ── En-tête / pied de page communs (refonte 2026) ─────────────────────
% #1 code fiche, #2 titre court, #3 couleur domaine
\newcommand{\ficheHeaders}[3]{%
  \fancyhead[L]{\small\sffamily{\bfseries\textcolor{#3}{#1}}\hspace{0.7em}\textcolor{colSlate}{#2}}%
  \fancyhead[R]{\small\sffamily\textcolor{colSlate}{CAPET SII IE\, --- \,Session 2026}}%
  \fancyfoot[C]{\small\sffamily\textcolor{colSlate}{\thepage\,/\,\pageref{LastPage}}}%
  \fancyhead[C]{}\fancyfoot[L]{}\fancyfoot[R]{}%
}

% ── Bloc en-tête de fiche (BLOC 0, refonte 2026) ──────────────────────
% #1 code  #2 titre  #3 couleur domaine  #4 domaine  #5 ligne présence  #6 sous-thèmes
\newcommand{\ficheEntete}[6]{%
  \begin{tcolorbox}[enhanced, breakable,
    frame hidden, boxrule=0pt, arc=2.5mm,
    interior style={left color=#3, right color=#3!75!black},
    left=9pt, right=9pt, top=7pt, bottom=7pt,
    before skip=0pt, after skip=6pt]
    {\color{white}\sffamily\bfseries\LARGE #1\quad #2}\\[4pt]
    {\color{#3!22!white}\sffamily\small \MakeUppercase{#4}\quad\textbar\quad #5}\\[1pt]
    {\color{#3!35!white}\sffamily\footnotesize Sous-thèmes : #6}
  \end{tcolorbox}%
}

% ── Mise en page ─────────────────────────────────────────────────────
\pagestyle{fancy}
\fancyhf{}
\renewcommand{\headrulewidth}{0.5pt}
\renewcommand{\footrulewidth}{0pt}
\AtBeginDocument{%
  \renewcommand{\headrule}{{\color{colRule}\hrule width\headwidth height 0.5pt}}%
}

\setlist[itemize]{leftmargin=1.2em, itemsep=1pt, topsep=2pt, parsep=0pt}
\setlist[enumerate]{leftmargin=1.2em, itemsep=3pt, topsep=2pt, parsep=0pt}

\setlength{\parskip}{2pt}
\setlength{\columnsep}{10pt}

\def\CAPETBEGINFILE{\begin{document}}
\def\CAPETENDFILE{\end{document}}
\else
\def\CAPETBEGINFILE{}
\def\CAPETENDFILE{}
\fi

\ficheHeaders{T11}{Machine à courant continu}{colElec}

\CAPETBEGINFILE

% -----------------------------------------------------------------------
%  BLOC EN-TETE
% -----------------------------------------------------------------------

\ficheEntete{T11}{Machine à courant continu (MCC)}{colElec}%
  {électrotechnique - Machines électriques}%
  {\textbf{Présence :} \textbf{10/13} sujets (2013-2026) - thème \emph{très fréquent}}%
  {équation électrique (FCEM, résistance d'induit) - couple électromagnétique - bilan de puissance - quadrants de fonctionnement - hacheur 1Q/4Q - démarrage rhéostatique}

% -----------------------------------------------------------------------
%  1. RAPPELS DE COURS
% -----------------------------------------------------------------------

\begin{tcolorbox}[
  ficheSection=colElec,
  title={1. Rappels de cours},
]

% -- 1.1 equations fondamentales ---------------------------------------
\textbf{1.1 Équations fondamentales (excitation séparée, flux $\phi$ constant)}

\smallskip
\begin{center}
\renewcommand{\arraystretch}{1.5}
\begin{tabular}{lll}
\toprule
\textbf{Grandeur} & \textbf{Expression} & \textbf{Convention} \\
\midrule
FCEM & $E = k\phi\Omega$ & $[k\phi] = \mathrm{V\!\cdot\!s\!\cdot\!rad^{-1}} = \mathrm{N\!\cdot\!m\!\cdot\!A^{-1}}$ \\[4pt]
Équation électrique (moteur) & $U = E + R_a I_a$ & $U > E$, courant entrant \\[4pt]
Équation électrique (générateur) & $U = E - R_a I_a$ & $E > U$, courant sortant \\[4pt]
Couple électromagnétique & $C_{em} = k\phi I_a$ & N$\cdot$m \\[4pt]
Vitesse & $\Omega = \dfrac{U - R_a I_a}{k\phi}$ & rad$\cdot$s$^{-1}$ \\
\bottomrule
\end{tabular}
\end{center}

\medskip
\hrule\vspace{4pt}

% -- 1.2 Bilan de puissance --------------------------------------------
\textbf{1.2 Bilan de puissance}

\smallskip
\[
  P_{abs} = U \cdot I_a = \underbrace{R_a I_a^2}_{P_{Joule}} + \underbrace{E \cdot I_a}_{P_{meca}}
  \qquad;\qquad
  P_{meca} = C_{em} \cdot \Omega
\]
Couple utile (en sortie d'arbre) : $C_u = C_{em} - C_{pertes}$ (pertes mécaniques, frottements).\\
Rendement : $\eta = P_{meca} / P_{abs} = E I_a / (U I_a)$.

\medskip
\hrule\vspace{4pt}

% -- 1.3 Quadrants de fonctionnement ----------------------------------
\textbf{1.3 Quadrants de fonctionnement}

\smallskip
\begin{center}
\renewcommand{\arraystretch}{1.25}
\begin{tabular}{llll}
\toprule
\textbf{Quadrant} & \textbf{Couple $C$} & \textbf{Vitesse $\Omega$} & \textbf{Mode} \\
\midrule
I  & $> 0$ & $> 0$ & Moteur direct \\
II & $< 0$ & $> 0$ & Générateur (frein récupératif) \\
III & $< 0$ & $< 0$ & Moteur inverse \\
IV & $> 0$ & $< 0$ & Générateur inverse \\
\bottomrule
\end{tabular}
\end{center}

\smallskip
\begin{center}
\begin{tikzpicture}[scale=0.95, font=\small]
  \draw[->] (-3.2,0) -- (3.2,0) node[right] {$C_{em}$};
  \draw[->] (0,-2.4) -- (0,2.4) node[above] {$\Omega$};
  % Famille de caracteristiques couple-vitesse (pente negative : caracteristique tombante)
  % Omega = U/(k*phi) - Ra/(k*phi)^2 * C_em : ordonnee a l'origine croissante avec U
  \draw[thick, colElec] (-2.6,2.25) -- (2.6,1.35);
  \draw[thick, colElec!55] (-2.6,1.45) -- (2.6,0.55);
  \draw[thick, colElec!30] (-2.6,0.65) -- (2.6,-0.25);
  \node[colElec, right] at (2.6,1.35) {\scriptsize $U_1$};
  \node[colElec!55, right] at (2.6,0.55) {\scriptsize $U_2 < U_1$};
  \node[colElec!30, right] at (2.6,-0.25) {\scriptsize $U_3 < U_2$};
  % Etiquettes des quadrants
  \node[font=\footnotesize] at (2.9,2.1) {I};
  \node[font=\footnotesize] at (-2.9,2.1) {II};
  \node[font=\footnotesize] at (-2.9,-2.1) {III};
  \node[font=\footnotesize] at (2.9,-2.1) {IV};
\end{tikzpicture}
\end{center}

\medskip
\hrule\vspace{4pt}

% -- 1.4 Hacheur de commande -------------------------------------------
\textbf{1.4 Hacheur de commande de la MCC}

\smallskip
\begin{center}
\renewcommand{\arraystretch}{1.3}
\begin{tabular}{lll}
\toprule
\textbf{Hacheur} & \textbf{Tension moyenne} & \textbf{Quadrants} \\
\midrule
1 quadrant (série) & $\langle U \rangle = \alpha \cdot E_{dc}$ & I seulement \\[4pt]
4 quadrants (pont en H bipolaire) & $\langle U \rangle = (2\alpha-1) \cdot E_{dc}$ & I, II, III, IV \\[4pt]
Démarrage - courant d'appel & $I_{demarr} = U/R_a$ ( $\Omega = 0$) & Limiter par rhéostats ou hacheur \\
\bottomrule
\end{tabular}
\end{center}

\end{tcolorbox}

% -----------------------------------------------------------------------
%  2. FORMULAIRE SYNTHETIQUE
% -----------------------------------------------------------------------

\begin{tcolorbox}[
  ficheSection=colElec,
  title={2. Formulaire synthétique},
]
\small
\begin{multicols}{2}
\renewcommand{\arraystretch}{1.45}
\begin{tabular}{@{}p{4.5cm}ll@{}}
\toprule
\textbf{Formule} & \textbf{Usage} & \textbf{U.} \\
\midrule
\multicolumn{3}{@{}l}{\footnotesize\textit{Équations électriques}} \\
$U = E + R_a I_a$ & Moteur & V \\
$U = E - R_a I_a$ & Générateur & V \\
$E = k\phi\Omega$ & FCEM & V \\
$C_{em} = k\phi I_a$ & Couple em & N$\cdot$m \\
$\Omega = (U - R_a I_a)/(k\phi)$ & Vitesse & rad/s \\
\midrule
\multicolumn{3}{@{}l}{\footnotesize\textit{Bilan de puissance}} \\
$P_{abs} = U I_a$ & Électrique & W \\
$P_J = R_a I_a^2$ & Pertes Joule & W \\
$P_{meca} = E I_a = C_{em}\Omega$ & Mécanique & W \\
\bottomrule
\end{tabular}

\columnbreak

\renewcommand{\arraystretch}{1.45}
\begin{tabular}{@{}p{4.5cm}ll@{}}
\toprule
\textbf{Formule} & \textbf{Usage} & \textbf{U.} \\
\midrule
\multicolumn{3}{@{}l}{\footnotesize\textit{Commande par hacheur}} \\
$\langle U \rangle = \alpha E_{dc}$ & Hacheur 1Q & V \\
$\langle U \rangle = (2\alpha-1)E_{dc}$ & Pont en H & V \\
$\alpha = \frac{1}{2}(1 + \langle U \rangle/E_{dc})$ & Rapport cyc. 4Q & - \\
\midrule
\multicolumn{3}{@{}l}{\footnotesize\textit{Couple nominal}} \\
$\Omega_{nom} = 2\pi N_{nom}/60$ & Conversion & rad/s \\
$C_{nom} = P_{nom}/\Omega_{nom}$ & Couple nom. & N$\cdot$m \\
$I_{demarr} = U/R_a$ & Courant à l'arrêt & A \\
\midrule
\multicolumn{3}{@{}l}{\footnotesize\textit{Données typiques}} \\
$E_{dc} = 207\,\mathrm{V}$ (PD2 + 230\,V) & Bus DC & V \\
$E_{dc} = 72\,\mathrm{V}$ (Urbanloop) & Bus DC & V \\
\bottomrule
\end{tabular}
\end{multicols}
\normalsize
\end{tcolorbox}

% -----------------------------------------------------------------------
%  3. QUESTIONS TYPES
% -----------------------------------------------------------------------

\begin{tcolorbox}[
  ficheSection=colElec,
  title={3. Questions types - extraites des sujets réels},
]

% -- QT1 -------------------------------------------------------------

\begin{tcolorbox}[questiontype,
  title={QT1 - Démarrage rhéostatique - courant constant (rame VMI) \hfill\normalfont\footnotesize E2 sujet 2019 (Q7-9)}
]
\textbf{Énoncé.}
4 moteurs CC en série sont alimentés par $V_{alim} = 750\,\mathrm{V}$. Lors du démarrage, le courant est maintenu à $I = 380\,\mathrm{A}$ constant par des rhéostats. La tension d'un moteur vaut $U_{moteur} = k_2 V_{train}^2$ avec $k_2 = 0{,}1\,\mathrm{V\!\cdot\!s^2\!\cdot\!m^{-2}}$. Donner la loi de commande $R_{th}(V_{train})$.

\tcbline

\textbf{Corrigé.}

\textbf{Étape 1 — Mise en équation.} Avec 4 moteurs en série et $R_a \approx 0$ :
\[
  V_{alim} = 4 U_{moteur} + R_{th} \cdot I
  \quad\Rightarrow\quad
  R_{th} = \frac{V_{alim} - 4 k_2 V_{train}^2}{I}
  = \frac{750 - 4 \times 0{,}1 \times V_{train}^2}{380}
\]

\textbf{Étape 2 — Interprétation physique.} La résistance \textbf{diminue} quand le train accélère (la FCEM $E = k_2 V_{train}^2$ augmente et prend en charge la tension). L'énergie dissipée dans les rhéostats correspond aux pertes par effet Joule lors du démarrage (toute l'énergie stockée dans les inductances).

\textbf{Conclusion.} La loi $R_{th}(V_{train})$ décroissante garantit un courant de démarrage constant à $380\,\mathrm{A}$ sur toute la plage de vitesse : elle respecte la contrainte de couple constant imposée par le cahier des charges de traction.
\end{tcolorbox}

\vspace{4pt}

% -- QT2 -------------------------------------------------------------

\begin{tcolorbox}[questiontype,
  title={QT2 - Hacheur 4 quadrants - rapport cyclique nominal \hfill\normalfont\footnotesize E2 appliquée 2022 (banc d'essai automobile, Q35-37)}
]
\textbf{Énoncé.}
Un frein à courant de Foucault est alimenté par un hacheur 4 quadrants (pont en H). $E_{dc} = 207\,\mathrm{V}$ (après redresseur PD2). Tension nominale du frein : $\langle u_f \rangle = 192\,\mathrm{V}$.
\begin{enumerate}[label=(\alph*)]
  \item Justifier le choix d'un hacheur 4 quadrants.
  \item Calculer $\alpha$ en commande bipolaire.
\end{enumerate}

\tcbline

\textbf{Corrigé.}

\textbf{Étape 1 — Justification (a).} Pour démagnétiser le frein, il faut inverser le courant (donc inverser la tension) : une réversibilité en tension et en courant est nécessaire, ce qui impose un \textbf{hacheur 4 quadrants} (pont en H).

\textbf{Étape 2 — Calcul de $\alpha$ (b).} Pont en H bipolaire : $\langle u_f \rangle = (2\alpha - 1) \cdot E_{dc}$
\[
  \alpha = \frac{1}{2}\left(1 + \frac{\langle u_f \rangle}{E_{dc}}\right)
  = \frac{1}{2}\left(1 + \frac{192}{207}\right)
  = \frac{1}{2}(1 + 0{,}927) \approx \mathbf{0{,}964}
\]

\textbf{Conclusion.} Le rapport cyclique $\alpha \approx 0{,}964$, proche de 1, confirme un fonctionnement en délestage quasi permanent côté positif : cohérent avec un frein fonctionnant majoritairement en régime établi et sollicité ponctuellement en inversion.
\end{tcolorbox}

\vspace{4pt}

% -- QT3 -------------------------------------------------------------

\begin{tcolorbox}[questiontype,
  title={QT3 - Couple nominal et vérification de dimensionnement \hfill\normalfont\footnotesize E2 sujet 2019 (Q28)}
]
\textbf{Énoncé.}
Les moteurs Alsthom ont : $P_{nom} = 96\,\mathrm{kW}$, $N_{nom} = 1035\,\mathrm{tr/min}$. Calculer le couple nominal. Sachant que le couple thermique équivalent sur le trajet est $C_{th} \approx 557\,\mathrm{N\!\cdot\!m}$, conclure sur le dimensionnement.

\tcbline

\textbf{Corrigé.}

\textbf{Étape 1 — Vitesse angulaire nominale.}
\[
  \Omega_{nom} = \frac{2\pi \times 1035}{60} = 108{,}4\,\mathrm{rad\!\cdot\!s^{-1}}
\]

\textbf{Étape 2 — Couple nominal.}
\[
  C_{nom} = \frac{P_{nom}}{\Omega_{nom}} = \frac{96\,000}{108{,}4} \approx \mathbf{885{,}6\,\mathrm{N\!\cdot\!m}}
\]

\textbf{Conclusion.} $C_{th} = 557\,\mathrm{N\!\cdot\!m} < C_{nom} = 885{,}6\,\mathrm{N\!\cdot\!m}$ : les moteurs sont \textbf{correctement dimensionnés} vis-à-vis du couple thermique équivalent du trajet.\checkmark
\end{tcolorbox}

\end{tcolorbox}

% -----------------------------------------------------------------------
%  4. ATTENTES DU JURY
% -----------------------------------------------------------------------

\begin{tcolorbox}[
  jury,
  title={4. Attentes du jury - Erreurs récurrentes et points de vigilance},
]

\begin{enumerate}[label=\textbf{\arabic*.}, leftmargin=*, itemsep=5pt]

  \item \textbf{Convention moteur vs générateur.}
  En moteur : $U > E$, $U = E + R_a I_a$. En générateur : $E > U$, $E = U + R_a I_a$.
  La confusion est explicitement sanctionnée par le jury 2024.

  \item \textbf{Au démarrage, $E = 0$.}
  À $\Omega = 0$, la FCEM est nulle : $I_{demarr} = U/R_a$ est très élevé.
  C'est pourquoi on insère des rhéostats de démarrage ou on utilise un hacheur (démarrage progressif).

  \item \textbf{Couple électromagnétique $\neq$ couple utile.}
  $C_{em} = k\phi I_a$ est le couple développé par les forces de Laplace.
  Le couple utile en sortie d'arbre est $C_u = C_{em} - C_{pertes\,mecaniques}$.

  \item \textbf{Unités du gain $k\phi$.}
  $[k\phi] = \mathrm{V\!\cdot\!s\!\cdot\!rad^{-1}} = \mathrm{N\!\cdot\!m\!\cdot\!A^{-1}}$ : le même coefficient apparaît dans $E = k\phi\Omega$ et dans $C_{em} = k\phi I_a$.

  \item \textbf{Hacheur 4Q obligatoire pour la réversibilité.}
  Un hacheur 1Q (tension toujours positive, courant toujours positif) ne permet pas le freinage récupératif ou l'inversion du sens de rotation. Justifier le choix du nombre de quadrants.

  \item \textbf{Bilan de puissance en 3 niveaux.}
  $P_{abs} = P_{Joule} + P_{meca}$. Présenter les trois termes avant tout calcul de rendement.

\end{enumerate}


\end{tcolorbox}

% -----------------------------------------------------------------------
%  PIED - PRESENCE DANS LES SUJETS
% -----------------------------------------------------------------------

\begin{tcolorbox}[
  ficheSection=colPresence,
  title={Présence dans les sujets analysés (2013-2026)},
]
\begin{center}
\footnotesize\renewcommand{\arraystretch}{1.15}
\begin{tabular}{lccccccccccccc}
\toprule
\textbf{Année} & 13 & 14 & 15 & 16 & 17 & 18 & 19 & 20 & 21 & 22 & 23 & 24 & 25 \\
\midrule
E1 (commune)   & - & $\checkmark$ & - & - & $\checkmark$ & - & - & $\checkmark$ & - & $\checkmark$ & - & - & $\checkmark$ \\
E2 discipl. IE & $\checkmark$ & - & $\checkmark$ & - & - & $\checkmark$ & $\checkmark$ & - & - & $\checkmark$ & - & $\checkmark$ & - \\
\bottomrule
\end{tabular}
\end{center}
\smallskip
\footnotesize\textbf{Sous-thèmes fréquents :}
équation électrique $U = E + R_a I_a$ (7/13) - couple électromagnétique (6/13) - bilan de puissance (5/13) - hacheur de commande 1Q/4Q (5/13) - démarrage rhéostatique (3/13)

\medskip
\textbf{Ordres de grandeur (repères rapides).}
\begin{itemize}[leftmargin=1.3em]
  \item Réseau BT FR : tension phase-neutre typique \SI{230}{\volt}, tension phase-phase typique \SI{400}{\volt}, fréquence \SI{50}{\hertz}.
  \item Courants d'étude fréquents en sujet : de quelques \si{\ampere} à quelques dizaines de \si{\ampere}.
  \item En conversion d'énergie, un rendement attendu reste généralement entre 0{,}8 et 0{,}98 selon la technologie.
\end{itemize}

\medskip
\textbf{Synthèse express (à refaire sans support).}
\begin{enumerate}[leftmargin=1.4em]
  \item Poser les hypothèses du modèle et les conventions de signe.
  \item Écrire 2 à 4 relations clés avec unités.
  \item Réaliser un calcul numérique complet et vérifier l'ordre de grandeur.
  \item Justifier le résultat par un argument technique de niveau CAPET.
\end{enumerate}

\end{tcolorbox}

\CAPETENDFILE
